\begin{mistake}{2026-04-12}{39}

\begin{problemcontent}
设 \(A,B\) 均为 \(n\) 阶反对称矩阵，且 \(AB = BA\)，则下列结论不正确的是（\quad）

(A) \(A+B\) 是反对称矩阵

(B) \(AB\) 是对称矩阵

(C) \(A^*+B^*\) 是反对称矩阵

(D) \(2A+3B\) 是反对称矩阵
\end{problemcontent}

\begin{solutioncontent}
由已知，\(A, B\) 为反对称矩阵，即 \(A^T = -A\)，\(B^T = -B\)，且 \(AB = BA\)。

对于 (A)：\((A+B)^T = A^T + B^T = -A - B = -(A+B)\)，故为反对称矩阵，正确。

对于 (B)：\((AB)^T = B^T A^T = (-B)(-A) = BA = AB\)，故为对称矩阵，正确。

对于 (D)：\((2A+3B)^T = 2A^T + 3B^T = 2(-A) + 3(-B) = -(2A+3B)\)，故为反对称矩阵，正确。

对于 (C)：
\[
\begin{aligned}
(A^*+B^*)^T &= (A^*)^T + (B^*)^T \\
&= (A^T)^* + (B^T)^* \\
&= (-A)^* + (-B)^* \\
&= (-1)^{n-1}A^* + (-1)^{n-1}B^* \\
&= (-1)^{n-1}(A^*+B^*)
\end{aligned}
\]
当 \(n\) 为奇数时，\((-1)^{n-1}=1\)，\(A^*+B^*\) 为对称矩阵；当 \(n\) 为偶数时，\((-1)^{n-1}=-1\)，\(A^*+B^*\) 为反对称矩阵。因 \(n\) 的奇偶性未知，故 (C) 的结论不正确。

故选 (C)。
\end{solutioncontent}

\mistakeknowledge{
\begin{itemize}
 \item \textbf{核心公式/定理}：反对称矩阵 \(A^T=-A\)；转置运算 \((AB)^T=B^TA^T\)；数乘伴随 \((kA)^*=k^{n-1}A^*\)。
 \item \textbf{易错点}：想当然地认为 \((-A)^* = -A^*\)。伴随矩阵元素为 \(n-1\) 阶代数余子式，提取负号需提 \(n-1\) 次。
 \item \textbf{关联知识}：伴随矩阵与转置可交换，即 \((A^*)^T=(A^T)^*\)；反对称矩阵的主对角线元素必然全为 \(0\)。
\end{itemize}}

\end{mistake}
